% !PRINTED 260
\weberpage{260}

% Running head "Siebenter Abschnitt." excluded as page furniture.

%<<CONT
Sind $\pi, \nu, \varrho$ die Anzahl der positiven, negativen, verschwindenden Quadrate, in die sich $\varphi$ zerlegen lässt, so sind $\pi, \nu$ die
Anzahlen der positiven und negativen Quadrate, in die sich die
durch (7) bestimmte Form $\Phi$ zerlegen lässt, und die Bestimmung der Zahlen $\pi, \nu$ braucht also nur noch für letztere Function, deren Determinante von Null verschieden ist, durchgeführt
zu werden.

\begin{center}
\S.\,82.

\gesperrt{Quadratische Formen mit nicht verschwindender\\
Determinante.}
\end{center}

Bei der Untersuchung der Formen $\varphi(x_1, x_2 \ldots x_m)$ mit nicht
verschwindender Determinante machen wir von einem Determinantensatz Gebrauch, den wir am Schluss des \S.\,28 bewiesen
haben, und an den wir hier erinnern wollen.

Ist $A$ eine Determinante von $m$ Reihen, und sind $A_i^{(k)}$ ihre
ersten, $A_{i,i'}^{k,k'}$ ihre zweiten Unterdeterminanten, so ist
\weq{1}{A_1^{(1)} A_i^{(i)} - A_1^{(i)} A_i^{(1)} = A A_{1,i}^{1,i}.}
Ist $A$ eine von Null verschiedene \gesperrt{symmetrische Determinante}, so ist $A_1^{(i)} = A_i^{(1)}$ und wir können die Formel (1) so
schreiben:
\weq{2}{A_1^{(1)} A_i^{(i)} - A_1^{(i)2} = A A_{1,i}^{1,i}.}
Darin kann $i$ jeden der Indices $2, 3 \ldots m$ bedeuten. Wenn
$A_1^{(1)} = 0$ und $A_{1,i}^{1,i} = 0$ ist, so folgt hieraus, dass auch $A_1^{(i)} = 0$
sein muss, und daraus schliessen wir, dass nicht zugleich
\[
A_1^{(1)}, \; A_{1,2}^{1,2}, \; A_{1,3}^{1,3} \ldots A_{1,m}^{1,m}
\]
verschwinden können, da sonst auch
\[
A_1^{(2)}, \; A_1^{(3)} \ldots A_1^{(m)},
\]
also auch, gegen die Voraussetzung, $A$ verschwinden würde (\S.\,22).

Wenden wir dies auf die jetzt von Null verschieden angenommene Determinante $R = R_m$ unserer Function $\varphi$ an, so
folgt, dass zwar die erste Haupt-Unterdeterminante $R_{m-1}$, dann
aber nicht alle zweiten Haupt-Unterdeterminanten $R_{m-2}$ verschwinden können. Dieselbe Schlussweise lässt sich anwenden,
wenn wir $R$ durch ein nicht verschwindendes $R_{m-1}, R_{m-2} \ldots$
ersetzen, und wir gelangen also zu folgendem wichtigen Satz:
%CONT>>
